% GENERATED FILE. Edit canonical lessons, then run npm run generate:books.
% canonical-source-hash: fnv1a64:36ccbbef20066641
% canonical-lessons: FR-C280-journal-televise, FR-C280-presentateur, FR-C280-telespectateur, FR-C280-auditeur, FR-C280-micro

\chapter{Le journal télévisé}
\label{ch:fr-le-journal-televise}
\hlchaptermodality{280}

\begin{chapteropening}
\textbf{By the end of this chapter:} \emph{I can talk about the news on television and radio: the TV news, the presenter, the viewers, the listeners and the microphone.}

Complete the last lesson of chapter 280: I can talk about the news on television and radio: the TV news, the presenter, the viewers, the listeners and the microphone.
\end{chapteropening}

\section[le journal télévisé]{le journal télévisé — the television news}

\label{lesson:FR-C280-journal-televise}



\begin{quote}
Before the new one: say the French for to come out, then the French for the publication.
\end{quote}

\subsection*{You'll want to know: le journal télévisé}
\textbf{le journal télévisé} — \textquotedblleft{}the television news\textquotedblright{}.

\begin{quote}
Nous regardons le journal télévisé à vingt heures.
\end{quote}

\begin{quote}
\emph{We watch the television news at eight in the evening.}
\end{quote}

Often \emph{le JT} in speech, or \emph{le journal de vingt heures}.

\begin{cousinweb}
\emph{journal} is inherited from Latin \emph{diurnālis} \textquotedblleft{}daily\textquotedblright{} (EtymDB); \emph{télévisé} is built on \emph{télévision}, from Greek \emph{tēle} \textquotedblleft{}far\textquotedblright{} and Latin \emph{vīsiō} \textquotedblleft{}sight\textquotedblright{}, from \emph{vidēre} \textquotedblleft{}to see\textquotedblright{}.
\end{cousinweb}

\begin{grammarlens}[title={What you've built}]
One more word for telling, explaining and describing at B1.
\end{grammarlens}

\subsection*{Your turn}
\begin{itemize}
\raggedright
  \item \emph{Say it:} \emph{le journal télévisé}
  \item \emph{Say it:} \emph{le journal télévisé}, once more
  \item \emph{Say it:} a sentence of your own with \emph{le journal télévisé}
  \item \emph{Recall:} say the French for to come out, then the French for the publication, then say \emph{le journal télévisé} again
\end{itemize}

\begin{tcolorbox}[breakable,colback=teal!4,colframe=teal!35!black,title={Before you move on}]
What does le journal télévisé mean? (\textquotedblleft{}The television news\textquotedblright{}.) Say it once more.
\end{tcolorbox}
\section[le présentateur]{le présentateur — the presenter, the newsreader}

\label{lesson:FR-C280-presentateur}



\begin{quote}
Before the new one: say the French for the publication, then the French for the television news.
\end{quote}

\subsection*{You'll want to know: le présentateur}
\textbf{le présentateur} — \textquotedblleft{}the presenter, the newsreader\textquotedblright{}.

\begin{quote}
Le présentateur a annoncé une bonne nouvelle.
\end{quote}

\begin{quote}
\emph{The presenter announced some good news.}
\end{quote}

Feminine: \emph{la présentatrice}.

\begin{cousinweb}
\emph{présentateur} is built on \emph{présenter} \textquotedblleft{}to present\textquotedblright{}, from Latin \emph{praesentāre} \textquotedblleft{}to show\textquotedblright{}, built on \emph{praesēns} \textquotedblleft{}present, at hand\textquotedblright{} (Lewis \& Short).
\end{cousinweb}

\begin{grammarlens}[title={What you've built}]
One more word for telling, explaining and describing at B1.
\end{grammarlens}

\subsection*{Your turn}
\begin{itemize}
\raggedright
  \item \emph{Say it:} \emph{le présentateur}
  \item \emph{Say it:} \emph{le présentateur}, once more
  \item \emph{Say it:} a sentence of your own with \emph{le présentateur}
  \item \emph{Recall:} say the French for the publication, then the French for the television news, then say \emph{le présentateur} again
\end{itemize}

\begin{tcolorbox}[breakable,colback=teal!4,colframe=teal!35!black,title={Before you move on}]
What does le présentateur mean? (\textquotedblleft{}The presenter, the newsreader\textquotedblright{}.) Say it once more.
\end{tcolorbox}
\section[le téléspectateur]{le téléspectateur — the viewer (of television)}

\label{lesson:FR-C280-telespectateur}



\begin{quote}
Before the new one: say the French for the television news, then the French for the presenter.
\end{quote}

\subsection*{You'll want to know: le téléspectateur}
\textbf{le téléspectateur} — \textquotedblleft{}the viewer (of television)\textquotedblright{}.

\begin{quote}
Des millions de téléspectateurs ont vu le match.
\end{quote}

\begin{quote}
\emph{Millions of viewers watched the match.}
\end{quote}

Feminine: \emph{la téléspectatrice}.

\begin{cousinweb}
\emph{téléspectateur} joins Greek \emph{tēle} \textquotedblleft{}far\textquotedblright{} and \emph{spectateur}, a learned borrowing from Latin \emph{spectātor} \textquotedblleft{}onlooker\textquotedblright{}, from \emph{spectāre} \textquotedblleft{}to watch\textquotedblright{} (Lewis \& Short).
\end{cousinweb}

\begin{grammarlens}[title={What you've built}]
One more word for telling, explaining and describing at B1.
\end{grammarlens}

\subsection*{Your turn}
\begin{itemize}
\raggedright
  \item \emph{Say it:} \emph{le téléspectateur}
  \item \emph{Say it:} \emph{le téléspectateur}, once more
  \item \emph{Say it:} a sentence of your own with \emph{le téléspectateur}
  \item \emph{Recall:} say the French for the television news, then the French for the presenter, then say \emph{le téléspectateur} again
\end{itemize}

\begin{tcolorbox}[breakable,colback=teal!4,colframe=teal!35!black,title={Before you move on}]
What does le téléspectateur mean? (\textquotedblleft{}The viewer (of television)\textquotedblright{}.) Say it once more.
\end{tcolorbox}
\section[l'auditeur]{l'auditeur — the listener}

\label{lesson:FR-C280-auditeur}



\begin{quote}
Before the new one: say the French for the presenter, then the French for the viewer.
\end{quote}

\subsection*{You'll want to know: l'auditeur}
\textbf{l'auditeur} — \textquotedblleft{}the listener\textquotedblright{}.

\begin{quote}
Les auditeurs peuvent appeler la radio.
\end{quote}

\begin{quote}
\emph{Listeners can phone in to the radio.}
\end{quote}

Feminine: \emph{l'auditrice}. In a firm, \emph{un auditeur} also checks the accounts.

\begin{cousinweb}
\emph{auditeur} is a learned borrowing from Latin \emph{audītor} \textquotedblleft{}hearer\textquotedblright{} (EtymDB), from \emph{audīre} \textquotedblleft{}to hear\textquotedblright{}.
\end{cousinweb}

\begin{grammarlens}[title={What you've built}]
One more word for telling, explaining and describing at B1.
\end{grammarlens}

\subsection*{Your turn}
\begin{itemize}
\raggedright
  \item \emph{Say it:} \emph{l'auditeur}
  \item \emph{Say it:} \emph{l'auditeur}, once more
  \item \emph{Say it:} a sentence of your own with \emph{l'auditeur}
  \item \emph{Recall:} say the French for the presenter, then the French for the viewer, then say \emph{l'auditeur} again
\end{itemize}

\begin{tcolorbox}[breakable,colback=teal!4,colframe=teal!35!black,title={Before you move on}]
What does l'auditeur mean? (\textquotedblleft{}The listener\textquotedblright{}.) Say it once more.
\end{tcolorbox}
\section[le micro]{le micro — the microphone, the mic}

\label{lesson:FR-C280-micro}



\begin{quote}
Before the new one: say the French for the viewer, then the French for the listener.
\end{quote}

\subsection*{You'll want to know: le micro}
\textbf{le micro} — \textquotedblleft{}the microphone, the mic\textquotedblright{}.

\begin{quote}
Le chanteur a pris le micro.
\end{quote}

\begin{quote}
\emph{The singer took the mic.}
\end{quote}

Short for \emph{le microphone}. \emph{Au micro} — speaking on air.

\begin{cousinweb}
\emph{micro} is short for \emph{microphone}, a learned word built from Greek \emph{mīkrós} \textquotedblleft{}small\textquotedblright{} and \emph{phōn\'{\={e}}} \textquotedblleft{}voice, sound\textquotedblright{} (EtymDB lists the short form without a parent).
\end{cousinweb}

\begin{grammarlens}[title={What you've built}]
One more word for telling, explaining and describing at B1.
\end{grammarlens}

\subsection*{Your turn}
\begin{itemize}
\raggedright
  \item \emph{Say it:} \emph{le micro}
  \item \emph{Say it:} \emph{le micro}, once more
  \item \emph{Say it:} a sentence of your own with \emph{le micro}
  \item \emph{Recall:} say the French for the viewer, then the French for the listener, then say \emph{le micro} again
\end{itemize}

\begin{tcolorbox}[breakable,colback=teal!4,colframe=teal!35!black,title={Before you move on}]
What does le micro mean? (\textquotedblleft{}The microphone, the mic\textquotedblright{}.) Say it once more.
\end{tcolorbox}
